% =====================================================================
% Kapitel 4: Datenmodell
% =====================================================================

\chapter{Datenmodell}

\label{ch:datenmodell}

Das Datenmodell ist das \emph{Fundament} des gesamten Projekts.
Es definiert, welche Daten gespeichert werden, wie sie
zusammenhängen und wie effizient wir sie abfragen können.

\section{Überblick: Die 14~Tabellen}

Tabelle~\ref{tab:tables} gibt einen schnellen Überblick über alle
14~Tabellen. Die vollständigen DDL-Definitionen finden sich in
Anhang~C.

\begin{table}[ht]
\centering
\small
\begin{tabular}{llll}
  \toprule
  \textbf{Tabelle} & \textbf{Zweck} & \textbf{PK} & \textbf{Wichtige FKs} \\
  \midrule
  \code{venues}              & Börsen/Handelsplätze       & \code{venue\_id}     & -- \\
  \code{vendors}            & Datenlieferanten           & \code{vendor\_id}    & -- \\
  \code{instruments}        & Asset-Master               & \code{instrument\_id} & \code{primary\_venue} \\
  \code{instrument\_identifiers} & Cross-Vendor-Mapping   & (\code{scheme}, \code{value}) & \code{instrument\_id} \\
  \code{trading\_calendars} & Handelstage je Venue       & (\code{venue\_id}, \code{date}) & \code{venue\_id} \\
  \code{prices\_raw}        & Roh-Preise (vendor-original) & (\code{vendor\_id}, \code{vendor\_symbol}, \code{date}) & \code{vendor\_id} \\
  \code{prices\_daily}      & OHLCV (harmonisiert)      & (\code{instrument\_id}, \code{date}) & \code{instrument\_id}, \code{vendor\_id} \\
  \code{corporate\_actions} & Splits, Dividenden         & (\code{instrument\_id}, \code{ex\_date}, \code{action\_type}) & \code{instrument\_id} \\
  \code{fx\_rates\_daily}   & Cross-Currency-FX          & (\code{ccy\_from}, \code{ccy\_to}, \code{date}) & \code{vendor\_id} \\
  \code{bond\_yields}        & Bond-Yield-Curve           & (\code{instrument\_id}, \code{date}) & \code{instrument\_id} \\
  \code{macro\_series}      & Makro-Stammdaten           & \code{series\_id}    & -- \\
  \code{macro\_observations}& Makro-Beobachtungen        & (\code{series\_id}, \code{date}) & \code{series\_id} \\
  \code{dq\_runs}            & DQ-Lauf-Metadaten          & \code{run\_id}       & -- \\
  \code{dq\_findings}        & DQ-Befunde                 & (\code{run\_id}, \code{check\_name}) & \code{run\_id} \\
  \bottomrule
\end{tabular}
\caption{Alle 14~Tabellen des \code{deephold\_db}-Schemas}
\label{tab:tables}
\end{table}

\section{ER-Diagramm}

Abbildung~\ref{fig:er} zeigt das vollständige ER-Diagramm.

\begin{figure}[ht]
\centering
\resizebox{\textwidth}{!}{%
\begin{tikzpicture}[
  table/.style={
    rectangle, draw, thick, rounded corners=2pt,
    minimum width=4cm, align=left, inner sep=4pt, font=\scriptsize,
  },
  title/.style={font=\bfseries, fill=blue!20, anchor=west, inner sep=2pt},
  col/.style={font=\ttfamily},
  pk/.style={col, fill=yellow!20},
  fk/.style={col, fill=green!20},
  arrow/.style={-Latex, thick},
]
  % instruments
  \node[table] (instruments) {
    \begin{tabular}{l}
      \node[title]{instruments}; \\
      \node[pk]{PK instrument\_id}; \\
      \node[col]{asset\_class (TEXT, NN)}; \\
      \node[col]{name (TEXT, NN)}; \\
      \node[col]{currency (CHAR(3))}; \\
      \node[fk]{FK primary\_venue $\to$ venues}; \\
      \node[col]{is\_active (BOOL)}; \\
    \end{tabular}
  };

  % venues
  \node[table, left=2.5cm of instruments] (venues) {
    \begin{tabular}{l}
      \node[title]{venues}; \\
      \node[pk]{PK venue\_id}; \\
      \node[col]{code (TEXT, UQ)}; \\
      \node[col]{name (TEXT)}; \\
      \node[col]{timezone (TEXT)}; \\
      \node[col]{country (CHAR(2))}; \\
    \end{tabular}
  };

  % vendors
  \node[table, above=1.2cm of instruments] (vendors) {
    \begin{tabular}{l}
      \node[title]{vendors}; \\
      \node[pk]{PK vendor\_id}; \\
      \node[col]{code (TEXT, UQ)}; \\
      \node[col]{license (TEXT)}; \\
      \node[col]{homepage (TEXT)}; \\
    \end{tabular}
  };

  % instrument_identifiers
  \node[table, right=2.5cm of instruments] (identifiers) {
    \begin{tabular}{l}
      \node[title]{instrument\_identifiers}; \\
      \node[pk]{PK (scheme, value)}; \\
      \node[fk]{FK instrument\_id $\to$ instruments}; \\
      \node[col]{scheme (TEXT, NN)}; \\
      \node[col]{value (TEXT, NN)}; \\
    \end{tabular}
  };

  % prices_daily
  \node[table, below=1.2cm of instruments] (prices) {
    \begin{tabular}{l}
      \node[title]{prices\_daily}; \\
      \node[pk]{PK (instrument\_id, date)}; \\
      \node[fk]{FK instrument\_id $\to$ instruments}; \\
      \node[fk]{FK vendor\_id $\to$ vendors}; \\
      \node[col]{open, high, low, close, adj\_close}; \\
      \node[col]{volume (BIGINT)}; \\
    \end{tabular}
  };

  % prices_raw
  \node[table, left=2.5cm of prices] (pricesraw) {
    \begin{tabular}{l}
      \node[title]{prices\_raw}; \\
      \node[pk]{PK (vendor\_id, vendor\_symbol, date)}; \\
      \node[fk]{FK vendor\_id $\to$ vendors}; \\
      \node[col]{payload (JSONB, NN)}; \\
    \end{tabular}
  };

  % macro_series
  \node[table, right=2.5cm of prices] (macroseries) {
    \begin{tabular}{l}
      \node[title]{macro\_series}; \\
      \node[pk]{PK series\_id (TEXT)}; \\
      \node[col]{name, source, frequency, unit}; \\
    \end{tabular}
  };

  % macro_observations
  \node[table, below=1.2cm of prices] (macroobs) {
    \begin{tabular}{l}
      \node[title]{macro\_observations}; \\
      \node[pk]{PK (series\_id, date)}; \\
      \node[fk]{FK series\_id $\to$ macro\_series}; \\
      \node[col]{value (NUMERIC(18,6))}; \\
    \end{tabular}
  };

  % FX
  \node[table, right=2.5cm of macroobs] (fx) {
    \begin{tabular}{l}
      \node[title]{fx\_rates\_daily}; \\
      \node[pk]{PK (ccy\_from, ccy\_to, date)}; \\
      \node[col]{rate (NUMERIC(18,8))}; \\
      \node[fk]{FK vendor\_id $\to$ vendors}; \\
    \end{tabular}
  };

  % corporate_actions
  \node[table, left=2.5cm of macroobs] (corpactions) {
    \begin{tabular}{l}
      \node[title]{corporate\_actions}; \\
      \node[pk]{PK (instrument\_id, ex\_date, action\_type)}; \\
      \node[fk]{FK instrument\_id $\to$ instruments}; \\
      \node[col]{value (NUMERIC(18,8))}; \\
    \end{tabular}
  };

  % bond_yields
  \node[table, right=2.5cm of macroobs] (bondyields) {
    \begin{tabular}{l}
      \node[title]{bond\_yields}; \\
      \node[pk]{PK (instrument\_id, date)}; \\
      \node[fk]{FK instrument\_id $\to$ instruments}; \\
      \node[col]{yield, price, duration, rating}; \\
    \end{tabular}
  };

  % dq
  \node[table, below=1.5cm of macroobs] (dq) {
    \begin{tabular}{l}
      \node[title]{dq\_runs / dq\_findings}; \\
      \node[col]{Lauf-Metadaten + Befunde}; \\
    \end{tabular}
  };

  % Verbindungen
  \draw[arrow] (venues) -- (instruments);
  \draw[arrow] (vendors) -- (prices);
  \draw[arrow] (instruments) -- (identifiers);
  \draw[arrow] (instruments) -- (prices);
  \draw[arrow] (instruments) -- (corpactions);
  \draw[arrow] (instruments) -- (bondyields);
  \draw[arrow] (macroseries) -- (macroobs);
\end{tikzpicture}%
}
\caption{Entity-Relationship-Diagramm des \code{deephold\_db}-Schemas.
  \code{instruments} ist der zentrale Asset-Master; \code{vendors} ist
  die Quelle aller Preise. Makro- und Equity-Daten sind
  \emph{parallel} organisiert, nicht in einer Hierarchie.}
\label{fig:er}
\end{figure}

\section{Stammdaten-Tabellen}

\subsection{\code{venues}}

\emph{Börsen und Handelsplätze}. Beispiele: NASDAQ, NYSE, XETRA,
LSE, TYO. Jede Zeile beschreibt einen Handelsplatz eindeutig über
den Code (\code{NASDAQ}, \code{NYSE}, \code{XETRA} \ldots). Die
\code{timezone} ermöglicht später korrekte Datums-Berechnungen über
Zeitzonen hinweg.

\subsection{\code{vendors}}

\emph{Datenlieferanten} -- die Quellen, aus denen wir Preise
beziehen. Beispiele: \code{fred}, \code{ecb}, \code{yahoo}, \code{demo}.
Die \code{license}-Spalte speichert die Nutzungsbedingungen (z.\,B.\
\foreign{public domain}, \foreign{CC BY 4.0}, \foreign{Yahoo ToS --
privater Gebrauch}).

\subsection{\code{instruments}}

\emph{Der Asset-Master} -- eine zentrale Tabelle, in der jedes
gehandelte Wertpapier genau einmal vorkommt: eine Aktie, ein ETF,
ein Index, eine Anleihe. \code{asset\_class} ist ein
\emph{Diskriminator} (\code{equity}, \code{index}, \code{bond\_gov},
\ldots), der bestimmt, in welchen weiteren Tabellen das Instrument
referenziert wird.

\subsection{\code{instrument\_identifiers}}

\emph{Das Cross-Vendor-Mapping}. Ein Instrument kann unter
verschiedenen IDs bei verschiedenen Vendoren bekannt sein:

\begin{itemize}
  \item Yahoo Finance: \code{AAPL}
  \item ISIN: \code{US0378331005}
  \item Bloomberg: \code{AAPL:US}
  \item FRED: \code{DEXUSEU} (für EUR/USD-FX)
\end{itemize}

\code{instrument\_identifiers} bildet jede dieser IDs auf die
\emph{interne} \code{instrument\_id} ab. Der zusammengesetzte
Primärschlüssel (\code{scheme}, \code{value}) verhindert Duplikate.

\section{Preis-Tabellen}

\subsection{\code{prices\_raw} -- Das unveränderte Original}

\code{prices\_raw} speichert die \emph{exakten Roh-Payloads} der
Vendoren, ohne Interpretation. Jede Zeile enthält:

\begin{itemize}
  \item \code{vendor\_id} -- welcher Lieferant
  \item \code{vendor\_symbol} -- dessen nativers Ticker-Symbol
  \item \code{date} -- Handelstag
  \item \code{payload} -- die vollständige rohe Antwort als JSONB
  \item \code{ingested\_at} -- Zeitstempel des Pulls
\end{itemize}

\begin{merke}
\code{prices\_raw} ist der \emph{Append-Only Audit-Trail}. Wir
schreiben hinein, ändern aber nie. Selbst wenn die harmonisierte
Tabelle \code{prices\_daily} korrigiert werden muss, ist die
Original-Antwort noch da.
\end{merke}

\subsection{\code{prices\_daily} -- Die harmonisierte Zeitreihe}

\code{prices\_daily} ist die \emph{Arbeitstabelle} für alle
Preis-basierten Analysen. Sie ist auf \code{instrument\_id}
\emph{normalisiert} (nicht auf das Vendor-Symbol) -- dadurch
können wir Daten verschiedener Vendoren für dasselbe Instrument
zusammenführen und vergleichen.

\begin{lstlisting}[language=sql, caption={Auszug aus \code{prices\_daily}}]
\d prices_daily
\end{lstlisting}

\begin{lstlisting}[language=sql]
                                       Table "public.prices_daily"
       Column        |           Type           | Nullable |    Default
--------------------+--------------------------+----------+---------------
 instrument_id       | bigint                   | not null |
 date                | date                     | not null |
 open                | numeric(18,8)            |          |
 high                | numeric(18,8)            |          |
 low                 | numeric(18,8)            |          |
 close               | numeric(18,8)            |          |
 adjusted_close      | numeric(18,8)            |          |
 volume              | bigint                   |          |
 vendor_id           | integer                  |          |
Indexes:
    "prices_daily_pkey" PRIMARY KEY, btree (instrument_id, date)
    "ix_prices_daily_date" btree (date)
    "ix_prices_daily_instrument_id_date" btree (instrument_id, date)
Check constraints:
    "prices_daily_check" CHECK (instrument_id IS NOT NULL)
Foreign-key constraints:
    "prices_daily_instrument_id_fkey" FOREIGN KEY (instrument_id) REFERENCES instruments(instrument_id)
    "prices_daily_vendor_id_fkey" FOREIGN KEY (vendor_id) REFERENCES vendors(vendor_id)
\end{lstlisting}

Beachten Sie die Indizes: \code{ix\_prices\_daily\_date} beschleunigt
\emph{zeitbasierte} Queries (\code{WHERE date BETWEEN ...}), der
zusammengesetzte Index \code{ix\_prices\_daily\_instrument\_id\_date}
beschleunigt Queries nach \emph{einem Instrument über die Zeit}.

\section{Makro-Tabellen}

\subsection{\code{macro\_series} und \code{macro\_observations}}

Makro-Daten (Zinssätze, Inflation, Arbeitsmarkt) werden in zwei
\emph{getrennte} Tabellen aufgeteilt:

\begin{itemize}
  \item \code{macro\_series} -- Stammdaten: Name, Quelle, Frequenz,
        Einheit.
  \item \code{macro\_observations} -- die eigentlichen Werte, eine
        Zeile pro (Serie, Datum).
\end{itemize}

Diese Aufteilung ist ein klassisches
\emph{Star-Schema}: Die \code{series}-Tabelle ist die
\emph{Dimension}, \code{observations} ist die \emph{Fact}-Tabelle.
Sie ermöglicht, eine Serie nachträglich um Metadaten zu ergänzen
(z.\,B.\ um die Datenquelle zu wechseln), ohne alle Beobachtungen
anzufassen.

\section{FX- und Bond-Tabellen}

\subsection{\code{fx\_rates\_daily}}

Wechselkurse -- eine Zeile pro (\code{ccy\_from}, \code{ccy\_to},
\code{date}). Die \code{ccy\_from}/\code{ccy\_to}-Codierung
statt einer \code{rate}-Tabelle mit zwei Zeilen pro Paar spart
Speicher und macht Queries schneller.

\subsection{\code{bond\_yields}}

Bond-spezifische Daten -- die allgemeine \code{prices\_daily}-Tabelle
enthält keine Yield-Curve-Daten. Hier werden zusätzlich
\code{yield}, \code{price}, \code{duration} und \code{rating}
gespeichert.

\section{DQ-Tabellen}

\subsection{\code{dq\_runs} und \code{dq\_findings}}

\emph{Data-Quality}-Tabellen. \code{dq\_runs} protokolliert jeden
Lauf eines DQ-Checks (z.\,B.\ \foreign{OHLC-Konsistenz} oder
\foreign{no negative prices}), \code{dq\_findings} die einzelnen
Befunde. Diese Tabellen sind derzeit nur \emph{Skelette} -- die
konkreten Checks werden in einem späteren Kapitel erläutert.

\section{SQLAlchemy-Definitionen}

Das Schema wird in Python deklariert. Listing~\ref{lst:models}
zeigt einen Auszug.

\begin{lstlisting}[language=python, caption={Auszug aus \code{src/deephold\_db/db/models.py}}]
class Instrument(Base):
    __tablename__ = "instruments"
    instrument_id: Mapped[int] = mapped_column(BigInteger, primary_key=True, autoincrement=True)
    asset_class: Mapped[str] = mapped_column(Text, nullable=False)
    name: Mapped[str] = mapped_column(Text, nullable=False)
    currency: Mapped[str | None] = mapped_column(CHAR(3))
    primary_venue: Mapped[int | None] = mapped_column(Integer, ForeignKey("venues.venue_id"))
    is_active: Mapped[bool] = mapped_column(Boolean, server_default=func.true())
    created_at: Mapped[datetime] = mapped_column(DateTime(timezone=True), server_default=func.now())

class PricesDaily(Base):
    __tablename__ = "prices_daily"
    instrument_id: Mapped[int] = mapped_column(
        BigInteger, ForeignKey("instruments.instrument_id"), primary_key=True
    )
    date: Mapped[date] = mapped_column(Date, primary_key=True)
    open: Mapped[Decimal | None] = mapped_column(Numeric(18, 8))
    high: Mapped[Decimal | None] = mapped_column(Numeric(18, 8))
    low: Mapped[Decimal | None] = mapped_column(Numeric(18, 8))
    close: Mapped[Decimal | None] = mapped_column(Numeric(18, 8))
    adjusted_close: Mapped[Decimal | None] = mapped_column(Numeric(18, 8))
    volume: Mapped[int | None] = mapped_column(BigInteger)
    vendor_id: Mapped[int | None] = mapped_column(Integer, ForeignKey("vendors.vendor_id"))
    __table_args__ = (
        Index("ix_prices_daily_date", "date"),
        Index("ix_prices_daily_instrument_id_date", "instrument_id", "date"),
    )
\end{lstlisting}

Listing~\ref{lst:models} zeigt die zentrale
\emph{Declarative}-Syntax von SQLAlchemy~2.x. Die
\code{Mapped[]}-Annotation gibt Typ + Optionalität an, der
\code{mapped\_column()}-Aufruf definiert die
Datenbank-Repräsentation.

\section{Was als Nächstes?}

Mit dem Schema im Hinterkopf wenden wir uns den \emph{Vendor-Adaptern}
zu -- den Komponenten, die Daten aus FRED, ECB und Yahoo Finance
in dieses Schema einfüllen. Kapitel~\ref{ch:vendor}.
